\chapter{Stage 3: The Three-Tier Router}
\label{ch:router}

We now have everything the router needs. Stage 1 produced a readiness signal,
$\upre$; Stage 2 produced, for the nearest stored episode, a similarity score and
a stored trust score. The router is the small piece of logic that turns those
three numbers into a single decision: which of the three tiers should answer this
query. It is the hinge of the whole per-query pipeline, the point where the
self-organising compute ladder of \Cref{ch:bigpicture} actually does its
organising.

\begin{intuition}[The cheapest tier that will do]
The router's guiding rule is economy under safety. It wants the \emph{cheapest}
tier that can answer well, so it prefers serving a trusted memory directly, then
generating zero-shot, and only then paying for retrieval. But it will spend more
the instant the cheap options stop being safe. Think of it as a dispatcher with
three crews of rising cost: it sends the cheapest crew that is clearly competent
for the job, and escalates the moment competence is in doubt. The art is in
defining ``clearly competent,'' and that is what the router's thresholds encode.
\end{intuition}

\section{What the router sees and what it returns}
\label{sec:router-io}

The router is deliberately starved of raw data. It does not touch the model, the
memory index, or the retrieval corpus; it only reads scalars the upstream stages
already computed. Its inputs are the pre-routing confidence $\upre$, the cosine
similarity $s$ of the query to its nearest stored episode, and that episode's
stored trust score $\ustored$. From these it produces a tier index in
$\{1, 2, 3\}$ together with bookkeeping: the routing score it computed, a flag
recording whether the safety override fired, the identifier of the retrieved
episode, and, when consulted, the output of the retrieval-utility classifier of
\Cref{ch:ruc}.

\begin{intuition}[Two numbers about memory, asking different questions]
The two memory numbers are not redundant, and the router treats them differently.
Similarity asks ``how close is the nearest thing we have seen?'' Stored trust asks
``and was that thing actually any good?'' A query can score high on one and low on
the other. We have a near-identical match to a past episode whose answer we never
trusted; or a loosely related match to an episode we trust completely. Keeping the
two questions separate is what lets the router make the right call in each case,
as we will see in \Cref{sec:router-asymmetry}.
\end{intuition}

\section{Three mechanisms, in a fixed order}
\label{sec:router-mechanisms}

The router applies three mechanisms, and the order is not negotiable, because each
expresses a different constraint that must dominate the ones after it.

\begin{description}[leftmargin=0pt, labelsep=0.5em, style=nextline]
  \item[Mechanism 1, the safety override (checked first).] If $\upre$ is below the
  floor $\phisafe = 0.38$, the query is dispatched off the score path to the
  retrieval-augmented tier and the override flag is set, regardless of how strong
  the memory match is. This is the veto from \Cref{ch:upre}: low pre-generation
  readiness means the system must not answer without grounded retrieval.
  \item[Mechanism 2, the dispatch score.] If the override does not fire, the router
  fuses the two memory numbers into one score and compares it against a threshold
  ladder, choosing Tier 1, Tier 2, or a fall-through.
  \item[Mechanism 3, the retrieval-utility classifier (Stage 3b).] The two
  outcomes that would otherwise route unconditionally to retrieval, the safety
  override and the score fall-through, are intercepted by a learned classifier that
  may redirect the query to zero-shot generation instead. The classifier is the
  subject of \Cref{ch:ruc}; here we note only where it plugs in.
\end{description}

The dispatch score itself is a weighted blend of the two memory numbers,
\[
  \text{routing score} \;=\; \lambda \cdot s \;+\; (1-\lambda)\cdot \ustored,
  \qquad \lambda = 0.70,
\]
which leans on similarity (weight $0.70$) more than on stored trust (weight
$0.30$), because a close match is the stronger evidence that memory covers the
query. The score then meets a two-rung ladder: a query whose score reaches the
combined threshold $\theta_1 = 0.90$ is served directly from memory (Tier 1); a
query that misses that rung but whose \emph{similarity alone} exceeds
$\theta_{\text{sim}} = 0.75$ is generated zero-shot (Tier 2); anything else falls
through toward retrieval (Tier 3), subject to the classifier interception.

\begin{table}[htbp]
\centering
\caption[Dispatch rules in order of evaluation]{The router's dispatch rules,
evaluated top to bottom; the first matching rule wins. The classifier of
\Cref{ch:ruc} intercepts the two rules that default to retrieval and may redirect
them to the zero-shot tier.}
\label{tab:dispatch}
\small
\renewcommand{\arraystretch}{1.25}
\begin{tabularx}{\textwidth}{@{}c l >{\raggedright\arraybackslash}X c@{}}
\toprule
\textbf{Order} & \textbf{Condition} & \textbf{Meaning} & \textbf{Tier} \\
\midrule
1 & $\upre < 0.38$ & safety override: readiness too low to generate ungrounded
  & 3 (or 2 via RUC) \\
2 & $0.70\,s + 0.30\,\ustored \ge 0.90$ & strong, trusted memory match: serve it
  & 1 \\
3 & $s > 0.75$ & similar enough to be worth a fresh zero-shot answer & 2 \\
4 & otherwise & no cheap path earns the query: retrieve and ground
  & 3 (or 2 via RUC) \\
\bottomrule
\end{tabularx}
\end{table}

\begin{lstlisting}[caption={The dispatch logic, condensed from \texttt{caem/routing/router.py}.}, label={lst:route}]
# Mechanism 1 -- safety override, evaluated FIRST
if u_pre < safety_thr:                            # safety_thr = 0.38
    return consult_ruc(default_tier=3, entry_point="safety_veto")  # may redirect to T2

# Mechanism 2 -- dispatch score
score = 0.70 * similarity + 0.30 * u_stored_retrieved
if score >= 0.90:                                 # tier1_combined_threshold
    tier = 1                                      # serve stored answer directly
elif similarity > 0.75:                           # tier2_similarity_threshold
    tier = 2                                      # generate zero-shot
else:
    tier = consult_ruc(default_tier=3, entry_point="fall_through")  # may redirect to T2
\end{lstlisting}

\section{Why the two tests are asymmetric}
\label{sec:router-asymmetry}

The single most instructive feature of the ladder is that the two rungs test
\emph{different things}. The Tier 1 rung tests the combined score, which needs both
similarity and stored trust to be high. The Tier 2 rung tests similarity alone, and
ignores stored trust entirely. This asymmetry is not an oversight; it follows
directly from what each tier does with the matched episode.

\begin{precise}[Serve-the-answer versus borrow-the-neighbourhood]
Tier 1 \emph{serves the stored answer verbatim}. To do that safely, the match must
be both close (so the stored answer is actually relevant) \emph{and} trusted (so the
stored answer is actually good). Both conditions are necessary, which is exactly why
the rung is a combined score with weight on each. Tier 2 does something different: it
\emph{generates a fresh answer}, using the existence of a similar past episode only
as evidence that the query lies in a region the model has handled before. The stored
answer is never served, so its quality is irrelevant to the Tier 2 decision; all that
matters is that a similar question exists, which is why the Tier 2 rung looks at
similarity alone. The asymmetry is the difference between reusing an answer and
reusing a neighbourhood.
\end{precise}

A concrete consequence is the case that a single-threshold router would get wrong: a
query with a near-identical match to an episode whose stored answer was never trusted.
The combined score stays below $0.90$ because the trust term is low, so Tier 1 is
correctly refused, the system will not serve the bad stored answer. But similarity is
high, comfortably above $0.75$, so the query routes to Tier 2 and the model generates a
fresh answer for a question it has clearly seen before. The router declines to reuse the
\emph{answer} while still exploiting the \emph{familiarity}.

\begin{figure}[htbp]
\centering
\begin{tikzpicture}[scale=6.0, font=\small\sffamily]
  % region fills
  \fill[cbScenario!14] (0,0) rectangle (0.75,1);                                   % Tier 3
  \fill[cbIntuition!16] (0.75,0) -- (1,0) -- (1,0.6667) -- (0.857,1) -- (0.75,1) -- cycle; % Tier 2
  \fill[cbExample!24] (0.857,1) -- (1,1) -- (1,0.6667) -- cycle;                    % Tier 1
  % boundaries
  \draw[dashed, cbInk!55, line width=0.7pt] (0.75,0) -- (0.75,1);
  \draw[cbExample!55!black, line width=0.9pt] (0.857,1) -- (1,0.6667);
  % axes
  \draw[->, cbInk!65, line width=0.8pt] (0,0) -- (1.06,0)
    node[right, font=\footnotesize] {similarity $s$};
  \draw[->, cbInk!65, line width=0.8pt] (0,0) -- (0,1.06)
    node[above, font=\footnotesize] {stored trust $\ustored$};
  % ticks
  \foreach \x in {0.25,0.5,0.75,1} \draw[cbInk!40] (\x,0) -- (\x,-0.02)
    node[below, font=\scriptsize] {\x};
  \foreach \y in {0.25,0.5,0.75,1} \draw[cbInk!40] (0,\y) -- (-0.02,\y)
    node[left, font=\scriptsize] {\y};
  % region labels
  \node[font=\footnotesize\bfseries, text=cbScenario!55!black, align=center]
    at (0.37,0.52) {Tier 3\\{\footnotesize retrieve (or RUC $\to$ T2)}};
  \node[font=\footnotesize\bfseries, text=cbIntuition!60!black, align=center]
    at (0.86,0.30) {Tier 2\\{\footnotesize zero-shot}};
  \node[font=\footnotesize\bfseries, text=cbExample!45!black, align=center, rotate=0]
    at (0.93,0.87) {Tier 1};
  % boundary annotation
  \node[font=\scriptsize, text=cbExample!45!black, anchor=west]
    at (0.45,0.92) {$0.70\,s + 0.30\,\ustored = 0.90$};
  \draw[cbExample!45!black, line width=0.4pt] (0.88,0.91) -- (0.905,0.86);
\end{tikzpicture}
\caption[The routing map]{The dispatch score and the similarity rung carve the
(similarity, stored-trust) plane into three regions. Tier 1 (green) is the top-right
corner where the combined score clears $0.90$: a close \emph{and} trusted match. Tier
2 (blue) is the rest of the high-similarity strip ($s > 0.75$): close enough to
generate fresh, even when the stored answer is not trusted. Tier 3 (teal) is
everything to the left: no cheap path is earned, so the query retrieves, subject to
the classifier of \Cref{ch:ruc}. The safety override of Mechanism 1 sits outside this
plane: when $\upre < 0.38$ it pre-empts the map entirely and forces the right edge. (Schematic.)}
\label{fig:routing-map}
\end{figure}

\section{Why three mechanisms instead of one threshold}
\label{sec:router-three}

The router's design separates three questions that a naive dispatcher would conflate.
The dispatch score captures \emph{whether memory covers the query}. The safety floor
captures \emph{whether the model is ready to generate at all}. The classifier of
\Cref{ch:ruc} captures \emph{whether retrieval will actually help on this particular
query}. Each is a genuinely different question, and answering them with one number
would force bad trades.

\begin{intuition}[What a single threshold would get wrong]
Suppose we collapsed everything into one score and one cutoff. Then a query with a
strong memory match but low readiness would be served from the cheap tier, because the
match dominates the number, exactly the confident-but-wrong case the safety floor
exists to prevent. Or a query with weak memory coverage would be sent to retrieval even
when retrieval is known to hurt it, because nothing in the single number asks whether
retrieval helps. By keeping the three questions on three mechanisms, the router serves
the cheap tier only when coverage and readiness \emph{both} say it is safe, while the
classifier governs the choice between zero-shot and retrieval on the queries that are
left.
\end{intuition}

\begin{defence}[If an examiner asks: ``how is this different from Adaptive-RAG?'']
Prior three-tier routers such as Adaptive-RAG~\cite{jeong2024adaptive} gate their tiers
with a single query-complexity classifier: the query's predicted difficulty alone
decides the tier. \caem{}'s router differs in two ways. It dispatches on \emph{per-entry
stored quality over a shared, evolving memory}, not on a static property of the query,
so the same query routes differently as the memory learns. And it separates the three
concerns above into three mechanisms evaluated in order, rather than folding them into
one difficulty score. The retrieval-utility classifier of \Cref{ch:ruc} is closest in
spirit to the Adaptive-RAG idea, but it sits \emph{downstream} of the memory-coverage
and readiness mechanisms, refining only the queries those mechanisms could not place on
a cheap tier.
\end{defence}

\section{The router at work}
\label{sec:router-scenarios}

Five readings cover the router's whole behaviour. The first is the cold start we have
already met; the rest are the four corners of \Cref{fig:routing-map} plus the override.

\begin{scenario}[Five queries, five routes]
\begin{enumerate}[leftmargin=1.5em, itemsep=2pt]
  \item \textbf{Cold start (empty memory).} No episode exists, so similarity and stored
  trust are both zero, the score is zero, and similarity is below $0.75$. The query falls
  through to Tier 3 (or the classifier redirects it to Tier 2). This is the FEVER claim of
  \Cref{ch:bigpicture} on cycle zero.
  \item \textbf{Close and trusted.} A near-identical match to a high-trust episode pushes
  the combined score past $0.90$. Tier 1: the stored answer is served directly, no
  generation, no verification.
  \item \textbf{Close but not trusted.} A near-identical match to a low-trust episode: the
  score stays under $0.90$ because the trust term drags it down, but similarity clears
  $0.75$. Tier 2: generate a fresh answer rather than serve the distrusted one.
  \item \textbf{Loosely related.} Similarity below $0.75$ and the score well under $0.90$.
  The query falls through to Tier 3, where the classifier decides whether retrieval will
  help or whether zero-shot is the better bet.
  \item \textbf{Unready, regardless of match.} The pre-routing confidence is below $0.38$.
  The safety override fires first and pre-empts the entire map: the query goes to the
  grounded tier (or the classifier sends it to zero-shot) no matter how good the memory
  match looked. This is the TriviaQA question of \Cref{ch:bigpicture}.
\end{enumerate}
\end{scenario}

These five readings are a conceptual map of the router's behaviour, but the deployed
run shows how unevenly the tiers are actually used. The cheap memory tier (Tier 1) is
exercised on only a tiny fraction of queries: none at the cold start, and roughly one
in a thousand once memory has grown (two of fifteen hundred at the third cycle, both
returning the correct stored answer). Live dispatch is therefore overwhelmingly a
choice between zero-shot generation (about forty-five percent of queries) and grounded
retrieval (about fifty-five percent), with Tier 1 acting as a precision-first cache for
near-verbatim repeats rather than a high-traffic path. The five-scenario walkthrough
still names every route the router can take; the deployed shares simply tell us which
of those routes carry the traffic.

\begin{figure}[htbp]
\centering
\begin{tikzpicture}
\begin{axis}[
    ybar stacked,
    width=\textwidth,
    height=6.0cm,
    bar width=22pt,
    ymin=0, ymax=100,
    ylabel={share of queries (\%)},
    symbolic x coords={fact-check, trivia, commonsense, multi-step, misconception},
    xtick=data,
    x tick label style={font=\footnotesize, align=center},
    ytick={0,20,40,60,80,100},
    enlarge x limits=0.12,
    legend style={at={(0.5,-0.20)}, anchor=north, legend columns=2, font=\footnotesize},
    legend cell align={left},
    every axis plot/.append style={fill},
]
  % Tier 2 (zero-shot generation) shares, per benchmark
  \addplot+[cbIntuition!60] coordinates {
    (fact-check,29.3) (trivia,31.7) (commonsense,79.7) (multi-step,23.3) (misconception,59.7)
  };
  % Tier 3 (grounded retrieval) shares, per benchmark
  \addplot+[cbScenario!50] coordinates {
    (fact-check,70.7) (trivia,68.0) (commonsense,20.3) (multi-step,76.7) (misconception,40.0)
  };
  \legend{Tier 2 (zero-shot generation), Tier 3 (grounded retrieval)}
\end{axis}
\end{tikzpicture}
\caption[Deployed per-benchmark tier shares at the third cycle]{Actual deployed tier
shares at the third cycle of the full run, per benchmark. The cheap memory tier
(Tier 1) is a hairline at this scale and is omitted. The commonsense probe routes
mostly to zero-shot generation (Tier 2), while the fact-checking, trivia, multi-step
reasoning, and misconception probes lean on grounded retrieval (Tier 3). This is the
measured counterpart to the conceptual regions of \Cref{fig:routing-map}.}
\label{fig:tier-share}
\end{figure}

\begin{bythenumbers}[Router constants]
\begin{itemize}[leftmargin=1.3em, itemsep=1pt]
  \item Score blend: $\lambda = 0.70$ on similarity, $0.30$ on stored trust.
  \item Tier 1 rung: combined score $\ge \theta_1 = 0.90$ (close \emph{and} trusted).
  \item Tier 2 rung: similarity $> \theta_{\text{sim}} = 0.75$ (close enough to generate).
  \item Safety override: $\upre < \phisafe = 0.38$, checked before the score.
  \item Evaluation order is fixed: override, then Tier 1 rung, then Tier 2 rung, then fall-through.
\end{itemize}
\end{bythenumbers}

The router has now chosen a tier for every query, but for two of its outcomes it deferred
to a classifier we have only named. \Cref{ch:ruc} opens that box, Stage 3b, the
retrieval-utility classifier: the learned component that decides, for a fall-through
query, whether retrieval will actually help or quietly hurt.
